\section{Introdução}

A lista ligada, também conhecida por cons list, é um pilar da programação funcional e uma das estruturas de 
dados mais simples e elegantes. Como as listas ligadas são apenas células ligadas sequencialmente por ponteiros, 
operações como cons (adicionar um elemento ao início), car (aceder ao primeiro elemento)
e cdr (aceder à lista excluindo o primeiro elemento) apenas requerem tempo e espaço $\mathcal{O}(1)$. No entanto, aceder 
a uma célula arbitrária da lista exige percorrer até $n - 1$ ponteiros para uma lista de comprimento $n$.
\\

Para resolver este problema, Eugene W. Myers \cite{myers1983} propôs uma nova estrutura de dados,
a Applicative Random-Access Stack, que denominamos Myers List. 
Essa estrutura mantém as operações de cons, car e cdr em tempo constante, permitindo ainda reduzir a complexidade 
de lookup para uma constante de $\mathcal{O}(\log n)$, para ser mais preciso $3 \log(n) - 5$ pointers. Esta melhoria é possível
observar pela adição de um ponteiro extra, denominado jump, em cada célula.
\\

Trabalhos mais recentes,  como o de Edward Peters, Yong Qi Foo e Michael D. Adams \cite{peters2025}, 
aprimoraram a construção e o uso desse ponteiro jump, resultando nas variantes 
apresentadas no paper, Improved Myers List e Advanced Myers List. Essas estruturas reduziram ainda mais o número de travessias
necessárias no pior caso, obtendo uma complexidade de $2 \log(n + 1) - 3$ e $(1 + \frac{1}{\sigma}) \log(n) + \sigma + 9$ respetivamente, 
aproximando-se do limite teórico de $\log(n + 1) - 1$ (limite teórico da informação).
\\

Apesar dessas melhorias, a Advanced Myers List apresenta uma limitação, a perda da funcionalidade de cons em tempo constante, 
funcionalidade que acreditamos ser de extrema importância. Além disso, durante todo o estudo do paper referenciado, os autores 
propõem algumas restrições, sendo uma delas a não adição de novos apontadores em cada célula. Isso motiva as duas perguntas 
centrais deste trabalho: \textbf{(i)} é possível recuperar a operação cons em tempo constante na Advanced Myers List, removendo 
a restrição do projeto, ou seja, permitindo a adição de ponteiros extras por célula enquanto se preserva a lógica e a estrutura 
apresentada no paper? \textbf{(ii)} removendo novamente a mesma restrição, é possível aplicar as ideias centrais propostas na 
Advanced Myers List à Improved Myers List, de modo a melhorar o seu desempenho de lookup e mantendo o cons constante?
\\

Neste relatório, propomos duas novas estruturas de dados, as suas respetivas implementações e bibliotecas, que respondem a estas
duas perguntas: a Dense Advanced Myers List (Secção \ref{subsec:dense_list}) e a Improved Applicative Random-Access Stack with Sigma region (IARASS) (Secção \ref{subsec:iarass}). 
Na secção \ref{sec:resultados} apresentamos os resultados experimentais de desempenho e de uso de memória. Na secção 5 discutimos as implementações destes
resultados face às perguntas colocadas.